\documentclass[11pt,a4paper]{article}
\usepackage[margin=2.4cm]{geometry}
\usepackage{microtype}
\usepackage{booktabs}
\usepackage{tabularx}
\usepackage{longtable}
\usepackage{array}
\usepackage{amsmath}
\usepackage{fancyvrb}
\usepackage{zkopipeline}
\usepackage[hidelinks]{hyperref}
\usetikzlibrary{positioning,arrows.meta}

\sisetup{group-separator={,},group-minimum-digits=4}
\setlength{\parindent}{0pt}
\setlength{\parskip}{0.6\baselineskip}
\fvset{fontsize=\small,frame=single,framesep=6pt,rulecolor=\color{zkodark!40}}

% ---------------------------------------------------------------------
% Data: everything below is computed from these two files at compile time
% ---------------------------------------------------------------------
\DTLloaddb{sites}{data/sites.csv}
\DTLloaddb{dict}{data/dictionary.csv}

\title{From Block Data to Slides\\[4pt]
  \large A reproducible pipeline for spatial energy analysis\\
  with QGIS, Python, \LaTeX\ and TikZ}
\author{Jan Suchanek}
\date{November 2026}

\begin{document}
\maketitle

\begin{abstract}\noindent
This document describes a pipeline that turns public, block-level energy data
into maps, charts and in-text figures for reports and slide decks. It was
developed in 2024--2025 for an internal analysis of energy use at eleven
science and technology sites in Berlin. That study is not public; this
document therefore uses \emph{synthetic data} for five fictitious sites.
The document is its own minimal working example: every chart and every number
in the text is computed from the data files when the document is compiled.
\end{abstract}

\tableofcontents

% =====================================================================
\section{Why a pipeline?}
% =====================================================================

The original question was simple: how much energy do the sites consume today,
from which sources, and how much will they need once they are fully built?
Answering it meant combining several public data sets, each at a different
spatial level, for eleven very different sites, and presenting the result
repeatedly as data were corrected and extended.

Doing this by hand (copying numbers from spreadsheets into slides) fails in a
predictable way: the numbers in the text drift away from the numbers in the
charts. The pipeline avoids this by keeping a single source of truth.
Every figure that appears in the output is \emph{computed} from one data file
at compile time. When the data change, recompiling updates every chart and
every sentence.

Section~\ref{sec:lessons} lists what went wrong in the first version, and how
this version prevents it. Those lessons are the main reason this document
exists.

% =====================================================================
\section{Architecture}\label{sec:architecture}
% =====================================================================

\begin{figure}[ht]
\centering
\begin{tikzpicture}[
  font=\sffamily\small,
  box/.style={draw=zkodark,rounded corners,fill=zkodark!6,align=center,
              minimum width=2.6cm,minimum height=1.2cm},
  file/.style={draw=zkolight,fill=zkolight!8,align=center,minimum width=2.6cm,
               minimum height=0.8cm},
  arr/.style={-{Stealth[length=6pt]},thick,zkodark}]
  \node[box] (raw) {Public data\\\scriptsize Energieatlas, ALKIS,\\[-2pt]\scriptsize IHK, polygons};
  \node[box,right=0.7cm of raw] (qgis) {QGIS\\\scriptsize spatial join};
  \node[file,right=0.7cm of qgis] (blocks) {\ttfamily blocks.csv\\\scriptsize one row per block};
  \node[box,below=0.9cm of blocks] (prep) {\ttfamily prepare.py\\\scriptsize aggregate, extrapolate};
  \node[file,left=0.7cm of prep] (sites) {\ttfamily sites.csv\\\ttfamily dictionary.csv};
  \node[box,left=0.7cm of sites] (tex) {\LaTeX\ + \ttfamily zkopipeline\\\scriptsize draw, compute, format};
  \node[box,below=0.9cm of tex,fill=zkolight!15] (out) {Report, slides, posters};
  \draw[arr] (raw) -- (qgis);
  \draw[arr] (qgis) -- (blocks);
  \draw[arr] (blocks) -- (prep);
  \draw[arr] (prep) -- (sites);
  \draw[arr] (sites) -- (tex);
  \draw[arr] (tex) -- (out);
\end{tikzpicture}
\caption{The pipeline. Boxes are processing steps, pink frames are files.
In the first version, the step \texttt{prepare.py} was a spreadsheet.}
\label{fig:pipeline}
\end{figure}

\begin{table}[ht]
\centering\small
\begin{tabularx}{\linewidth}{@{}l l X@{}}
\toprule
Stage & Tool & Responsibility \\
\midrule
1 Raw data & public portals & Block-level consumption, buildings, businesses,
  site boundaries as GeoJSON \\
2 Spatial join & QGIS & Assign blocks to sites; export one row per block \\
3 Preparation & Python & Aggregate per site, extrapolate missing values,
  derive indicators; write the data dictionary \\
4 Presentation & \LaTeX, TikZ, datatool & Draw maps and charts; compute
  every number used in the text \\
5 Output & PDF & Report, slide deck, one poster per site \\
\bottomrule
\end{tabularx}
\caption{Stages and their responsibilities.}
\end{table}

Two rules hold the design together:
\begin{enumerate}
  \item \textbf{One direction.} Data flow from left to right only. Nothing is
        typed into a later stage that could be derived from an earlier one.
  \item \textbf{Machine format inside, local format outside.} All files use a
        decimal point and full precision. Rounding and the decimal comma (for
        German output) are applied only when a number is printed.
\end{enumerate}

% =====================================================================
\section{Stage 1 and 2: data and spatial join}
% =====================================================================

\subsection{Sources}

The Berlin \emph{Umweltatlas} and \emph{Energieatlas} publish electricity, gas
and district heating consumption for each statistical city block. Further
layers came from the property register (ALKIS: buildings, listed monuments),
the chamber of commerce (IHK: registered businesses with size classes) and the
state statistical office. The boundaries of the sites were available as
GeoJSON polygons.

Two properties of the block data shape everything that follows:
\begin{itemize}
  \item A block can \emph{have} consumption without a \emph{published} value.
        Values are withheld when a block has very few consumers (data
        protection), or when a supplier did not deliver data.
  \item Withheld blocks are not a random sample. Few consumers per block often
        means one large consumer, so withheld blocks tend to consume
        \emph{more} than average.
\end{itemize}

\subsection{Selection rule: within or intersecting?}

A site boundary rarely follows block boundaries. QGIS offers two natural
selections: blocks lying \emph{within} the polygon, and blocks that
\emph{intersect} it. Figure~\ref{fig:selection} shows the difference.

\begin{figure}[ht]
\centering
\begin{tikzpicture}[font=\sffamily\scriptsize,scale=0.75]
  % blocks: 0 = outside, 1 = intersects, 2 = within
  \foreach \row/\vals in {0/{0,1,1,1,0,0},1/{1,2,2,2,1,0},2/{1,2,2,2,2,1},
                          3/{0,1,2,2,2,1},4/{0,0,1,1,1,0}}{
    \foreach \v [count=\col] in \vals {
      \ifcase\v
        \draw[gray!40] (\col,-\row) rectangle ++(0.92,-0.92);
      \or
        \filldraw[draw=zkolight,fill=zkolight!25] (\col,-\row) rectangle ++(0.92,-0.92);
      \or
        \filldraw[draw=zkodark,fill=zkodark!45] (\col,-\row) rectangle ++(0.92,-0.92);
      \fi
    }}
  \draw[very thick,zkodark,rounded corners=6pt]
    (1.55,-0.5) -- (4.6,-0.3) -- (6.4,-1.6) -- (6.5,-3.2) -- (4.2,-4.6)
    -- (2.6,-4.3) -- (1.3,-2.9) -- cycle;
  \node[anchor=west,align=left] at (7.4,-1.2)
    {\tikz\filldraw[draw=zkodark,fill=zkodark!45] (0,0) rectangle (0.3,0.3); within\\[3pt]
     \tikz\filldraw[draw=zkolight,fill=zkolight!25] (0,0) rectangle (0.3,0.3); intersects\\[3pt]
     \tikz\draw[very thick,zkodark] (0,0.15) -- (0.3,0.15); site boundary};
\end{tikzpicture}
\caption{Schematic selection of blocks by a site boundary.}
\label{fig:selection}
\end{figure}

\emph{Within} underestimates (it drops blocks that are partly on the site);
\emph{intersects} overestimates (it adds blocks that are mostly outside).
The original study used both, compared the results, and redrew one boundary by
hand where neither rule was acceptable. In this example, \texttt{prepare.py}
uses both relations; the rule is a single constant (\texttt{INCLUDE}) so that
its effect can be tested.

The QGIS export is a plain CSV with one row per block: the site it belongs to,
the relation (\texttt{within} or \texttt{intersects}), area, number of
buildings, and for each energy carrier a status (\texttt{reported},
\texttt{suppressed}, \texttt{none}) and the value where reported.

% =====================================================================
\section{Stage 3: preparation}\label{sec:prepare}
% =====================================================================

The script \texttt{data/prepare.py} replaces what was a spreadsheet in the
first version. It has no dependencies beyond the Python standard library.

\subsection{Extrapolating withheld values}

For each site $s$ and carrier $c$, let $B_{s,c}$ be the blocks with
consumption and $R_{s,c}\subseteq B_{s,c}$ the blocks with a reported value
$e_b$. The estimate assumes that withheld blocks consume as much as the
average reported block:
\begin{equation}
  \hat E_{s,c} \;=\; \frac{1}{|R_{s,c}|}\sum_{b\in R_{s,c}} e_b \;\cdot\; |B_{s,c}|
\end{equation}
This estimator has two known weaknesses, both documented in the output:
\begin{itemize}
  \item Because withheld blocks tend to be larger consumers, $\hat E$ is
        biased \emph{downwards}.
  \item If no block of a site reports a value ($|R_{s,c}|=0$), the estimate is
        zero although consumption exists. In the synthetic data this happens
        for district heating at \emph{Old Airfield}: heat is used in
        \DTLforeach*{sites}{\i=id,\hb=heat-blocks}{\ifthenelse{\equal{\i}{c}}{\hb}{}}
        blocks, but none reports a value.
\end{itemize}

\subsection{Emissions}

CO$_2$ emissions use one emission factor per carrier, in tonnes per GWh of
supplied energy. The values are those of the state energy balance used in the
original project; they are constants at the top of the script.

\begin{center}\small
\begin{tabular}{@{}lr@{}}
\toprule
Carrier & t\,CO$_2$/GWh \\
\midrule
Electricity & 365.22 \\
Natural gas & 202.77 \\
District heating & 219.65 \\
\bottomrule
\end{tabular}
\end{center}

\subsection{Build-out factor}

To estimate demand at full build-out, the current building density $d_s$
(buildings per hectare) is compared with a target density $d^\ast = 10$, the
density of the most built-up site in the original study:
\begin{equation}
  g_s = \max\!\left(1,\ \frac{d^\ast}{d_s}\right), \qquad
  E_s^{\text{future}} = g_s \cdot E_s
\end{equation}
The factor is an upper bound: it counts buildings, not floor area, and it
assumes that new buildings consume like existing ones. The $\max(1,\cdot)$ is
essential; without it, sites that are already denser than the target would
\emph{shrink} in the projection (see Section~\ref{sec:lessons}).

\subsection{Output: data and dictionary}

The script writes two files. \texttt{sites.csv} has one row per site and one
column per quantity, sorted by total demand. \texttt{dictionary.csv} describes
every column. Both are read by \LaTeX; the table below is typeset directly
from the dictionary.

% =====================================================================
\section{The data dictionary}\label{sec:dict}
% =====================================================================

{\small
\begin{longtable}{@{}>{\ttfamily}l p{7.6cm} l l@{}}
\toprule
\normalfont Column & Description & Unit & Source \\
\midrule
\endhead
\DTLforeach*{dict}{\dcol=column,\ddesc=description,\dunit=unit,\dsrc=source}{%
  \DTLiffirstrow{}{\\}\dcol & \ddesc & \dunit & \dsrc}\\
\bottomrule
\end{longtable}}

In the first version the dictionary lived in two extra rows at the bottom of
the data file (description and source), which the drawing macros skipped with
a numeric test. Moving it to its own file keeps the data table rectangular and
purely numeric, which every tool can read without special cases.

% =====================================================================
\section{Stage 4: drawing with \texttt{zkopipeline}}\label{sec:draw}
% =====================================================================

The package needs only the data loaded with \texttt{datatool}:
\begin{Verbatim}
\usepackage{zkopipeline}
\DTLloaddb{sites}{data/sites.csv}
\end{Verbatim}

\subsection{Circle map}

\verb|\zkomap| places the sites on a ring, connected by thin lines (the
network motif of the original project). The \emph{area} of each circle is
proportional to the value, so the radius grows with its square root. The
centre shows sum and mean.

\begin{Verbatim}
\zkomap[unit=GWh, scale=0.12]{sites}{total-gwh}
\end{Verbatim}

\begin{center}
\zkomap[unit=GWh,scale=0.12]{sites}{total-gwh}
\end{center}

\subsection{Projection overlay}

With \texttt{future=<column>}, each value is multiplied by the factor in that
column and drawn as a dashed circle. Sites whose factor is 1 get no overlay.

\begin{Verbatim}
\zkomap[unit=GWh, scale=0.08, future=factor]{sites}{total-gwh}
\end{Verbatim}

\begin{center}
\zkomap[unit=GWh,scale=0.08,future=factor]{sites}{total-gwh}
\end{center}
% keep the totals of this map for the text (the next chart overwrites \zkoresult)
\let\sitesnow\zkoresult \let\sitesfuture\zkofutureresult

\subsection{Energy mix}

\verb|\zkopies| draws one pie per site with three shares, here electricity,
gas and district heating. The pie area is proportional to the total. The label
gives the three shares in per cent.

\begin{Verbatim}
\zkopies[scale=0.12]{sites}{el-gwh}{gas-gwh}{heat-gwh}
\end{Verbatim}

\begin{center}
\zkopies[scale=0.12]{sites}{el-gwh}{gas-gwh}{heat-gwh}
\end{center}

\subsection{Bar chart}

Circle areas are good for showing orders of magnitude and poor for precise
comparison. Where readers must compare values, \verb|\zkobars| is the better
choice. Bars appear in the order of the data, which \texttt{prepare.py} sorts
by total demand. With \texttt{color=<column>}, each bar takes its colour
from that column (a colour name per row), e.g.\ to match the colours of the same
categories in a map.

\begin{Verbatim}
\zkobars[unit=GWh, width=8]{sites}{total-gwh}
\end{Verbatim}

\begin{center}
\zkobars[unit=GWh,width=8]{sites}{total-gwh}
\end{center}

% =====================================================================
\section{Numbers in the text}\label{sec:numbers}
% =====================================================================

Every drawing macro leaves its results in global macros: \verb|\zkoresult|
holds the sum, \verb|\zkofutureresult| the projected sum. The text then
computes with these values instead of copying them.

\subsection{An example paragraph}

\zkototal{sites}{gas-gwh}{\gas}%
\zkototal{sites}{co2-t}{\co}%
The source of the following paragraph contains no typed number:

\begin{quote}
The five sites draw \zkonum{\sitesnow}~GWh of energy per year, of which natural
gas accounts for \zkoshare{\gas}{\sitesnow}. Their supply causes about
\zkonum[0]{\co/1000}~kt of CO$_2$. At full build-out, demand could rise to
\zkonum{\sitesfuture}~GWh, a change of \zkogrowth{\sitesnow}{\sitesfuture}, or
\zkoratio{\sitesnow}{\sitesfuture} times today's level.
\end{quote}

\begin{Verbatim}
\zkototal{sites}{gas-gwh}{\gas}  \zkototal{sites}{co2-t}{\co}
The five sites draw \zkonum{\sitesnow}~GWh of energy per year, of which
natural gas accounts for \zkoshare{\gas}{\sitesnow}. Their supply causes
about \zkonum[0]{\co/1000}~kt of CO$_2$. At full build-out, demand
could rise to \zkonum{\sitesfuture}~GWh, a change of
\zkogrowth{\sitesnow}{\sitesfuture}, or \zkoratio{\sitesnow}{\sitesfuture} times ...
\end{Verbatim}

The values \verb|\sitesnow| and \verb|\sitesfuture| were taken from the
projection map above, directly after drawing it:
\verb|\let\sitesnow\zkoresult \let\sitesfuture\zkofutureresult|.
They must be saved at once, because the next chart overwrites
\verb|\zkoresult|.

\subsection{Share, growth, ratio}

Three macros, three different questions. Keeping them apart by name prevents
the most expensive error of the first version.

\begin{center}\small
\begin{tabular}{@{}lll@{}}
\toprule
Macro & Computes & Example: from 2,145 to 11,026 \\
\midrule
\verb|\zkoshare{a}{b}|  & $100\cdot a/b$ & \zkoshare{11026}{2145} \quad(share; wrong here)\\
\verb|\zkogrowth{a}{b}| & $100\cdot (b-a)/a$ & \zkogrowth{2145}{11026} \quad(correct)\\
\verb|\zkoratio{a}{b}|  & $b/a$ & \zkoratio{2145}{11026} \\
\bottomrule
\end{tabular}
\end{center}

Each macro takes an optional first argument to store the unrounded result,
e.g.\ \verb|\zkogrowth[\g]{a}{b}|, and prints \emph{n/a} instead of dividing
by zero.

\subsection{Rounding and locale}

Values are rounded only when printed, by \texttt{siunitx}. The same value can
be shown in English or German notation by changing one setting:

\begin{center}\small
\begin{tabular}{@{}ll@{}}
\toprule
Setting & Output of \verb|\zkonum[2]{\sitesfuture}| \\
\midrule
English (this document) & \zkonum[2]{\sitesfuture} \\
German & {\sisetup{output-decimal-marker={,},group-separator={.}}\zkonum[2]{\sitesfuture}} \\
\bottomrule
\end{tabular}
\end{center}

\begin{Verbatim}
\sisetup{output-decimal-marker={,}, group-separator={.}}
\end{Verbatim}

% =====================================================================
\section{Lessons learned}\label{sec:lessons}
% =====================================================================

The first version produced a 121-slide deck, a report and eleven posters. A
review of that version found the following errors. None was a mistake in the
data; each was a weakness of the pipeline. The right column shows how this
version rules it out.

{\small
\begin{longtable}{@{}p{4.6cm}p{5.2cm}p{5.2cm}@{}}
\toprule
Symptom & Cause & Remedy in this version \\
\midrule
\endhead
A 5.14-fold increase was reported as ``growth of 514\,\%'' instead of 414\,\%. &
One macro computed $a/b$ and was used both for shares and for growth. &
Separate, explicitly named macros: \verb|\zkoshare|, \verb|\zkogrowth|,
\verb|\zkoratio|. \\
\addlinespace
Three percentages in the text did not match the charts. &
Numbers had been typed into the text from an earlier data version. &
No typed numbers: every value in the text is read or computed from the data
(Section~\ref{sec:numbers}). \\
\addlinespace
Projected growth shown as +144\,\% instead of +137.9\,\%. &
Values were rounded to 0.1\,TWh \emph{before} dividing. &
Full precision in all files and computations; rounding only on output. \\
\addlinespace
Nine dated copies of the data file. &
Versioning by file name; old numbers survived in old copies. &
One data file, generated by a script and versioned with Git. \\
\addlinespace
Derived columns could not be reproduced. &
The spreadsheet kept values only; the formulas survived as text descriptions. &
\texttt{prepare.py} computes every derived column; the dictionary documents it. \\
\addlinespace
Mixed decimal separators in the output; conversion code in every macro. &
Data stored with decimal comma, but TikZ and \texttt{fp} need a point. &
Decimal point in all files; \texttt{siunitx} formats the output. \\
\addlinespace
Dense sites shown as ``could grow by $-3\,\%$''; the projected total too low. &
Build-out factors below 1 were applied and summed. &
$g = \max(1, d^\ast/d)$ in \texttt{prepare.py} (Section~\ref{sec:prepare}). \\
\addlinespace
Electricity use (GWh) and PV capacity (MWp) both labelled ``GW''. &
One unit argument for a chart that combined two quantities. &
Each column has its own unit in the dictionary; one quantity per chart. \\
\addlinespace
``Division by zero'' for a site without data. &
The zero test compared strings: \texttt{0.00} is not \texttt{0}. &
Numerical tests only: \verb|\zkotest{a == 0}{...}{...}|. \\
\addlinespace
``Dimension too large'' with large values (e.g.\ tonnes of CO$_2$). &
Comparisons via \verb|\ifdim|, limited to about 16,383. &
All arithmetic and comparisons with \verb|\fpeval|. \\
\addlinespace
Errors with some \TeX\ installations only. &
\LaTeX\ macros (site names) inside a data file, read before they were defined. &
Data files contain plain text only. \\
\addlinespace
Duplicate labels, broken cross-references (``Folie ??''). &
Frames copied to create variants, labels copied with them. &
Labels derived from a single naming scheme; warnings treated as errors. \\
\bottomrule
\end{longtable}}

% =====================================================================
\section{Building}\label{sec:build}
% =====================================================================

\begin{Verbatim}
python3 data/make_example_data.py   # synthetic blocks.csv (skip with real data)
python3 data/prepare.py             # sites.csv and dictionary.csv
pdflatex doku.tex                   # twice, for the table of contents
\end{Verbatim}

Requirements: Python~3.8 or later (standard library only) and a current \TeX\
distribution with \texttt{tikz}, \texttt{datatool}, \texttt{siunitx}. Tested
with \TeX\ Live 2023 (\texttt{datatool} 2.32) and \texttt{pdflatex}.

To use real data, replace \texttt{blocks.csv} by the QGIS export (same
columns) and adjust \texttt{SITE\_NAMES} in \texttt{prepare.py}.

% =====================================================================
\appendix
\section{Macro reference}
% =====================================================================

{\small
\begin{tabularx}{\linewidth}{@{}l X@{}}
\toprule
Macro & Purpose \\
\midrule
\verb|\zkomap[opt]{db}{col}| & Circle map; stores \verb|\zkoresult|,
  with \texttt{future} also \verb|\zkofutureresult| \\
\verb|\zkopies[opt]{db}{c1}{c2}{c3}| & One pie per site with three shares \\
\verb|\zkobars[opt]{db}{col}| & Horizontal bars in data order \\
\verb|\zkototal{db}{col}{\cs}| & Column sum, full precision, into \verb|\cs| \\
\verb|\zkonum[digits]{expr}| & Evaluate and print a number, rounded \\
\verb|\zkoshare[\cs]{a}{b}| & $100\,a/b$ in per cent \\
\verb|\zkogrowth[\cs]{a}{b}| & $100\,(b-a)/a$ in per cent, with sign \\
\verb|\zkoratio[\cs]{a}{b}| & $b/a$ \\
\verb|\zkotest{expr}{true}{false}| & Numerical condition \\
\midrule
Options & \texttt{unit}, \texttt{scale} (cm per $\sqrt{\text{unit}}$),
  \texttt{ring} (radius), \texttt{future} (factor column),
  \texttt{width} (longest bar, cm), \texttt{mesh} (true/false),
  \texttt{color} (bars: column with a colour name per row) \\
\bottomrule
\end{tabularx}}

\end{document}
